To begin our formal description of the prior work,
we'll introduce some  notation.  
A dataset consists of $n$ \emph{examples}, 
or \emph{data points}
$\{\vct x_i \in \mathcal{X}, y_i \in \mathcal{Y} \}$,
each consisting of a feature vector $\vct x_i$ and a label $y_i$.  A supervised learning algorithm $f:\mathcal{X}^n \times \mathcal{Y}^n \rightarrow 
(\mathcal{X} \rightarrow [0,1])$ is a mapping from datasets to models.  The learning algorithm
produces a model 
$\hat{y}: \mathcal{X} \rightarrow \mathcal{Y}$,
which given a feature vector $\vct x_i$,
predicts the corresponding output $y_i$.
% The supervised learning algorithm $f$ 
% is itself a function,
% mapping from datasets to models  
% $f:(\mathcal{X}^n, \mathcal{Y}^n) \rightarrow 
% (\mathcal{X} \rightarrow [0,1])$.
In this discussion
we focus on binary classification
%the setting in which the label $y$ 
% takes values from the set 
($\mathcal{Y} = \{0,1\}$).

We consider probabilistic classifiers
which produce estimates $\hat{p}(\mathbf{x})$ of the conditional probability 
$\P(y=1 \mid \vct x)$
of the label given a feature vector $\mathbf{x}$. 
To make a prediction $\hat{y}(\vct x) \in \mathcal{Y}$
given an estimated probability $\hat{p}(\vct x)$ a threshold rule is used: $\hat{y}_i = 1 \text{ iff } \hat{p}_i > t$.
The optimal choice of the threshold $t$
depends on the performance metric being optimized.  In our theoretical analysis, we consider optimizing the \textit{immediate utility}
%, as introduced in 
\citep{corbett2017algorithmic}, of which classification accuracy (expected $0-1$ loss) is a special case.   We will define this metric more precisely in the next section.  


In formal descriptions of discrimination-aware ML,
a dataset  possesses a protected feature $z_i \in \mathcal{Z}$,
making each example a three-tuple $(\vct x_i, y_i, z_i)$.
The protected characteristic may be real-valued, like age, 
or categorical, like race or gender.  The goal of many methods in discrimination-aware ML is not only to maximize accuracy, but also to ensure some form of impact parity.
Following related work, 
we consider binary protected features 
that divide the set of examples 
into two groups $a$ and $b$.  Our analysis extends directly to settings with more than two groups.  


Of the various measures of impact disparity,
%that have been proposed, 
the two that are the most relevant here are the Calders-Verwer gap and the p-\% rule. At a given threshold $t$, let
\mbox{$q_z = \frac{1}{n_z}\sum_{i : z_i=z}  \mathbbm{1}(\hat p_i > t)$}, where $n_z = \sum_{i}^{n} \mathbbm{1}(z_i = z)$. 
The \textbf{Calders-Verwer (CV) gap},
$q_a - q_b$,
is the difference between the proportions 
assigned to the positive class 
in the advantaged group $a$ and the disadvantaged  group $b$ \citep{kamishima2011fairness}.
The p-\% rule 
is a related metric\cite{zafar2017fairness}. 
Classifiers satisfy the \textbf{p-\% rule} 
if $q_b / q_a \ge p/100$.  
% This metric is motivated by 
% fair employment practices \citep{biddle2006adverse}, 
% in which it is stated that cases where the 
% ratio $q_b/q_a$ is below $0.8$ are problematic.  


Many papers in discrimination-aware ML 
propose to optimize accuracy (or some other risk) 
subject to constraints on the resulting 
level of impact parity 
as assessed by some metric \citep{pedreshi2008discrimination,kamiran2010discrimination,dwork2017decoupled,bechavod2017learning,hardt2016equality,ritov2017conditional}.  
% Papers proposing DLPs \citep{pedreshi2008discrimination,kamiran2009classifying, zafar2017fairness} 
% take as a premise 
Use of DLPs presupposes
that using the protected feature $z$ 
as a model input is impermissible in this effort. 
% as it amounts to \emph{treatment disparity}.
Discarding protected features, however, 
does not guarantee impact parity \cite{dwork2012fairness}.  
% Instead of discarding the protected features, 
DLPs incorporate $z$ in the learning algorithm, 
but without making it an input to the classifier.  
Formally, a DLP is a mapping: 
$\mathcal{X}^n \times \mathcal{Y}^n \times \mathcal{Z}^n \rightarrow (\mathcal{X} \rightarrow \mathcal{Y})$.
% AS $z$ isn't an input 
% to the classifier model, 
By definition, DLPs achieve treatment parity.  
However, satisfying \emph{treatment parity} in this fashion 
may still violate \emph{disparate treatment}.
% As we show next (\S \ref{section:theory}), 
% Under some circumstances, 
% DLPs can indirectly realize any function 
% achievable through treatment disparity.  

\paragraph{Alternative approaches.} 
Researchers have proposed a number of other techniques for reconciling accuracy and impact parity. 
One approach consists of preprocessing
the training data 
to reduce the dependence between 
the resulting model predictions 
and the sensitive attribute \citep{kamiran2009classifying,kamiran2012data,
feldman2015certifying,adler2016auditing,
johndrow2017algorithm}. 
These methods differ in terms of 
which variables they affect
and the degree of independence achieved.  
\cite{kamiran2009classifying} 
proposed flipping negative labels 
of training examples form the protected group.
\cite{zemel2013learning} proposed 
learning representations (cluster assignments)
so that group membership 
cannot be inferred from cluster membership. 
\cite{feldman2015certifying} 
and \cite{johndrow2017algorithm} 
also construct representations
designed to be marginally independent from $Z$.  
% \citet{johndrow2017algorithm} 
% demonstrates 
% how to construct transformations 
% to ensure that the derived features 
% are jointly independent of $Z$, 
% and show that this produces distributional parity 
% of the resulting fitted model. 
% A second common approach 
% is to modify existing classification methods 
% either through post-hoc corrections 
% or in the training stage 
% to constrain the level of impact parity 
% in the resulting model. 
% \citet{kamishima2011fairness, goh2016satisfying, calders2010three,kamiran2010discrimination} 
% consider modifications to methods such as SVM, 
% logistic regression, Naive Bayes, and decision trees.  
% \citet{agarwal2017reductions} 
% show how impact parity constraints 
% can be framed as a cost-sensitive classification problem. 


